\documentclass[11pt,letterpaper]{article}
\usepackage{technical-user-manual}
\setDocTitle{GitHub Actions Workflow Engineering Reference}
\setDocSubtitle{Engineering technical reference}
\setDocVersion{1.0}
\setDocStatus{Technical education / proposed reference}
\setDocOwner{CI/CD Engineering}
\setDocDate{30 October 2026}
\setDocSummary{Supplementary engineering practice}
\begin{document}
\makedoctitle
\makedocsummary
\tableofcontents
\clearpage
\section{Purpose, scope and authority}
This is an original modern engineering reference. The 2007 Duvall, Matyas and Glover book informs foundational CI intent but does not prescribe GitHub-specific implementations or 2026 supply-chain standards.

\section{Event and execution model}
Understand workflow events, triggers, jobs, steps, dependencies, concurrency groups, artifacts and permissions. Treat pull requests from forks as low trust; avoid deploying their head code with privileged tokens. Keep examples outside active workflow directory until reviewed.

Verification: collect the configuration or runtime evidence supporting this control, identify an accountable owner, and record the exception path where the condition cannot presently be met.

\section{Least privilege}
Set default permissions to contents: read, elevate narrowly per job and disable persisted checkout credentials unless required. Code scanning uploads may need security-events: write. Cloud OIDC needs id-token: write only in a protected deployment job.

Verification: collect the configuration or runtime evidence supporting this control, identify an accountable owner, and record the exception path where the condition cannot presently be met.

\section{Integrity}
Production actions should be pinned by verified full commit SHA or organization-approved lockfile workflow. Dependabot/Renovate policy should drive updates. Runner, container images, third-party actions and dependency downloads are supply-chain dependencies.

Verification: collect the configuration or runtime evidence supporting this control, identify an accountable owner, and record the exception path where the condition cannot presently be met.

\section{Separation of duties}
Use dedicated jobs, immutable artifacts, protected GitHub Environments and audited approvals. Validate production authorization in the cloud trust policy, not solely through GitHub job if-statements. Named environment review rules depend on repository plan and settings.

Verification: collect the configuration or runtime evidence supporting this control, identify an accountable owner, and record the exception path where the condition cannot presently be met.

\section{Operational behavior}
Use cancellation only where safe; avoid canceling a deployment midway through irreversible state transitions. Use concurrency to serialize prod but preserve state. Distinguish artifact handoff from cross-trust cache restoration.

Verification: collect the configuration or runtime evidence supporting this control, identify an accountable owner, and record the exception path where the condition cannot presently be met.

\section{Reference documents}
Consult the GitHub workflow syntax, workflow security, OIDC, cache, artifact attestation and dependency review documentation in references/source-register.tex before adoption.

Verification: collect the configuration or runtime evidence supporting this control, identify an accountable owner, and record the exception path where the condition cannot presently be met.

\section{Adoption checklist}
\begin{itemize}
\item Identify applicable repositories, application owners and trust boundaries.
\item Define minimum pass criteria and measurable acceptance thresholds.
\item Review workflow permissions, external actions, outputs and retention.
\item Record exceptions with justification, owner and expiry.
\item Revalidate after action, runner or policy changes.
\end{itemize}
\end{document}
